%% PATRÓN · Secuencia de llamadas: una calle vertical por componente (línea clara de su color), cada
\providecolor{textochip}{named}{white} % por si la charla tiene un estilo/ anterior a 1.1.0
%% llamada un chip en su calle, flechas en orden. Original: charla01/ap/llamadas.tex.
\begin{tikzpicture}[
    ch/.style={fill=#1,text=textochip,rounded corners=1.5pt,inner sep=2.5pt,font=\tiny\bfseries,minimum width=1.9cm},
    flecha/.style={-{Stealth[length=4pt]},tinta,line width=0.7pt}]
  \node[font=\scriptsize\bfseries,text=serieA] at (0,0.55) {MOMENTUM};
  \node[font=\scriptsize\bfseries,text=serieC] at (2.5,0.55) {CONTINUITY};
  \draw[serieA!30!fondo,line width=0.6pt] (0,0.3) -- (0,-3.2); \draw[serieC!30!fondo,line width=0.6pt] (2.5,0.3) -- (2.5,-3.2);
  \node[ch=serieA] (a) at (0,0) {paso 1: $\boldsymbol u$};
  \node[ch=serieC] (b) at (2.5,-0.9) {paso 3: $z$};
  \node[ch=serieA] (c) at (0,-1.8) {paso 4: $\boldsymbol v$};
  \node[ch=serieC] (d) at (2.5,-2.7) {paso 9: $z$};
  \draw[flecha] (a) -- (b); \draw[flecha] (b) -- (c); \draw[flecha] (c) -- (d);
  \node[font=\scriptsize,align=center,anchor=north] at (1.25,-3.3) {1 iteración\\ $\Rightarrow$ 2 llamadas a cada solver};
\end{tikzpicture}
